\documentclass[border=15pt]{standalone}
\usepackage{tikz}
\usetikzlibrary{arrows.meta, positioning, calc, fit, backgrounds}
\usepackage{amsmath}
\usepackage{xcolor}

% Colors
\definecolor{episodecolor}{RGB}{70, 70, 70}
\definecolor{intervalcolor}{RGB}{31, 119, 180}
\definecolor{roundcolor}{RGB}{44, 160, 44}
\definecolor{elementcolor}{RGB}{214, 39, 40}
\definecolor{solvercolor}{RGB}{148, 103, 189}
\definecolor{calloutcolor}{RGB}{255, 127, 14}

\begin{document}
\begin{tikzpicture}[
    node distance=0.55cm,
    >={Stealth[length=3mm]},
    step/.style={rectangle, draw, rounded corners=3pt, minimum height=0.8cm, 
                 align=center, font=\normalsize},
    callout/.style={font=\normalsize\itshape, text=calloutcolor!85!black},
    looplabel/.style={font=\small\itshape, text=black!55},
]

% ============================================================
% STEP 1: Episode Start
% ============================================================
\node[step, fill=episodecolor!8, draw=episodecolor, minimum width=14.5cm] (start) 
    {\textbf{1}\quad Sample random IC from pool of 7 waveforms;\; initialize level-1 mesh (8 elements)};

% ============================================================
% STEP 2: Compute errors
% ============================================================
\node[step, fill=intervalcolor!8, draw=intervalcolor!70, minimum width=13cm, 
      below=0.9cm of start] (errors)
    {\textbf{2}\quad Compute error indicators $e_k$ for all active elements};

% ============================================================
% STEP 3: Set thresholds
% ============================================================
\node[step, fill=intervalcolor!8, draw=intervalcolor!70, minimum width=13cm,
      below=0.4cm of errors] (thresh) 
    {\textbf{3}\quad Set thresholds:\quad
     $e_{\max} = \alpha \cdot \|e\|_\infty$\;,\quad
     $e_{\min} = e_{\max}^{\,\beta}$
     \quad\textit{\small(fixed for all rounds)}};

% ============================================================
% STEP 4: Assemble queue
% ============================================================
\node[step, fill=roundcolor!8, draw=roundcolor!70, minimum width=11cm,
      below=0.8cm of thresh] (queue)
    {\textbf{4}\quad Assemble element queue};

% Queue callout (below)
\node[callout, below=0.05cm of queue] (qcall) 
    {$\longrightarrow$\; \textsf{Component 5: Queue Mechanics}};

% ============================================================
% STEP 5: Observe, mask, decide, execute
% ============================================================
\node[step, fill=elementcolor!8, draw=elementcolor!70, minimum width=9cm,
      below=0.65cm of qcall] (visit)
    {\textbf{5}\quad Observe $\to$ Mask $\to$ Decide $\to$ Execute};

% Observation callout (right of step 5)
\node[callout, right=0.4cm of visit.east, anchor=west] (obscall) 
    {$\longrightarrow$\; \textsf{Component 4: Observation Space}};

% ============================================================
% STEP 6: Balance + local reward
% ============================================================
\node[step, fill=elementcolor!8, draw=elementcolor!70, minimum width=9cm,
      below=0.4cm of visit] (balance)
    {\textbf{6}\quad Enforce 2:1 balance;\; compute local reward};

% Reward callout (right of step 6)
\node[callout, right=0.4cm of balance.east, anchor=west] (rewcall) 
    {$\longrightarrow$\; \textsf{Component 6: Dual Reward (local)}};

% ============================================================
% Element loop arrow
% ============================================================
\draw[->, thick, elementcolor!60] (balance.south) -- ++(0, -0.4) -| 
    ([xshift=-5.2cm]visit.west) -- (visit.west)
    node[looplabel, pos=0.35, below, xshift=0.3cm] {next element in queue};

% ============================================================
% STEP 7: Round complete
% ============================================================
\node[step, fill=roundcolor!8, draw=roundcolor!70, minimum width=11cm,
      below=1.2cm of balance] (rebuilt)
    {\textbf{7}\quad Round complete\;---\;rebuild queue from updated mesh};

% ============================================================
% Round loop arrow
% ============================================================
\draw[->, thick, roundcolor!60] (rebuilt.south) -- ++(0, -0.4) -| 
    ([xshift=-6.2cm]queue.west) -- (queue.west)
    node[looplabel, pos=0.35, below, xshift=0.6cm] {next round\; (\texttt{max\_level} rounds total)};

% ============================================================
% STEP 8: Solver advance
% ============================================================
\node[step, fill=solvercolor!8, draw=solvercolor!70, minimum width=13cm,
      below=1.2cm of rebuilt] (solver)
    {\textbf{8}\quad Solver advances PDE by $T$;\quad track\; $\hat{e}_k = \displaystyle\max_{t \in [t_\tau,\, t_\tau + T]} e_k(t)$};

% ============================================================
% STEP 9: Global reward
% ============================================================
\node[step, fill=solvercolor!8, draw=solvercolor!70, minimum width=13cm,
      below=0.4cm of solver] (global)
    {\textbf{9}\quad Compute global retrospective reward};

% Global reward callout (below)
\node[callout, below=0.05cm of global] (gcall)
    {$\longrightarrow$\; \textsf{Component 6: Dual Reward (global)}};

% ============================================================
% Remesh interval loop arrow
% ============================================================
\draw[->, thick, intervalcolor!60] ([yshift=-0.15cm]gcall.south) -- ++(0, -0.3) -| 
    ([xshift=-7.4cm]errors.west) -- (errors.west)
    node[looplabel, pos=0.35, below, xshift=1.0cm] {next remesh interval\; ($N_{\text{remesh}}$ total)};

% ============================================================
% STEP 10: Episode End
% ============================================================
\node[step, fill=episodecolor!8, draw=episodecolor, minimum width=14.5cm,
      below=1.3cm of gcall] (end)
    {\textbf{10}\quad Episode End\quad 
     (${\approx}\; N_{\text{remesh}} \times \texttt{max\_level} \times n_{\text{active}} \approx 180$ RL steps)};

% ============================================================
% Main flow arrows
% ============================================================
\draw[->, thick] (start) -- (errors);
\draw[->, thick] (errors) -- (thresh);
\draw[->, thick] (thresh) -- (queue);
\draw[->, thick] (queue) -- (qcall);
\draw[->, thick] (qcall) -- (visit);
\draw[->, thick] (visit) -- (balance);
\draw[->, thick] (rebuilt) -- (solver);
\draw[->, thick] (solver) -- (global);
\draw[->, thick] (global) -- (gcall);

% ============================================================
% NESTED BOXES (background layer)
% ============================================================
\begin{scope}[on background layer]

  % --- Element Visit box (innermost) ---
  \node[draw=elementcolor!45, fill=elementcolor!3, rounded corners=5pt, 
        inner xsep=10pt, inner ysep=8pt,
        fit=(visit)(balance)(obscall)(rewcall),
        label={[font=\small\bfseries, text=elementcolor!70]above left:Element Visit}] 
        (elembox) {};

  % --- Adaptation Round box (middle) ---
  \node[draw=roundcolor!45, fill=roundcolor!3, rounded corners=6pt,
        inner xsep=12pt, inner ysep=10pt,
        fit=(queue)(qcall)(elembox)(rebuilt),
        label={[font=\small\bfseries, text=roundcolor!70]above left:Adaptation Round}]
        (roundbox) {};

  % --- Remesh Interval box (outer) ---
  \node[draw=intervalcolor!45, fill=intervalcolor!3, rounded corners=7pt,
        inner xsep=14pt, inner ysep=12pt,
        fit=(errors)(thresh)(roundbox)(solver)(global)(gcall),
        label={[font=\small\bfseries, text=intervalcolor!70]above left:Remesh Interval}]
        (intbox) {};

\end{scope}

% ============================================================
% Footer note
% ============================================================
\node[below=0.6cm of end, font=\small\itshape, text=black!60] (footer)
    {Orange labels mark pivot points\;---\;pause the walkthrough and explain the subsystem on a separate board.};

\end{tikzpicture}
\end{document}
